\chapter*{Executive Summary}
\addcontentsline{toc}{chapter}{Executive Summary}

The Netherlands faces a persistent housing shortage while construction is constrained by permitting delays, labour scarcity, financing costs, and environmental limits. Because historical extrapolation cannot fully represent these interacting constraints, this study applies Cooke's Classical Model for Structured Expert Judgment to forecast four indicators for 2027: completed dwellings, residential building permits, the average transaction price in Rotterdam, and free-sector rent growth.

Seven domain experts supplied 5th, 50th, and 95th percentiles for ten calibration variables and four target variables. Their statistical accuracy and informativeness were evaluated against the known calibration realizations. The optimized significance threshold was $\alpha^*=\OptimalAlpha$. Five experts received non-zero performance weights; two were excluded because their calibration scores were below the threshold. The optimized Decision Maker achieved a calibration score of \OptimizedDMCalibration{} and a mean information score of \OptimizedDMInformation{}, outperforming the equal-weight Decision Maker on the combined score.

The optimized median forecasts are 74,847 completed dwellings, 67,915 residential building permits, a Rotterdam transaction price of \euro669,298, and free-sector rent growth of 4.96\%. The corresponding 90\% intervals remain wide, showing that the panel agrees more strongly on the direction of housing pressure than on its exact magnitude. In particular, the median construction forecast remains well below the government's 100,000-dwelling annual ambition. The results should be interpreted as a transparent synthesis of expert uncertainty rather than a deterministic prediction, especially given the small, network-based panel and the limited set of ten calibration variables.
